\documentclass[12pt,a4paper,german]{report}
\usepackage[utf8]{inputenc}
\usepackage[ngerman]{babel}
\usepackage[left=2cm,right=2cm,top=2cm,bottom=2cm]{geometry}
\usepackage{amsmath}
\usepackage{amssymb}
\usepackage{listings}
\usepackage{xcolor}
\usepackage{hyperref}
\usepackage{fancyhdr}
\usepackage{array}
\usepackage{booktabs}
\usepackage{graphicx}
\usepackage{parskip}
\usepackage{pifont}

% Farben für Code-Blöcke
\definecolor{codegreen}{rgb}{0,0.6,0}
\definecolor{codegray}{rgb}{0.5,0.5,0.5}
\definecolor{codepurple}{rgb}{0.58,0,0.82}
\definecolor{backcolour}{rgb}{0.95,0.95,0.92}

% Listings-Konfiguration
\lstset{
    backgroundcolor=\color{backcolour},
    commentstyle=\color{codegreen},
    keywordstyle=\color{codepurple},
    numberstyle=\tiny\color{codegray},
    stringstyle=\color{codegreen},
    basicstyle=\ttfamily\small,
    breakatwhitespace=false,
    breaklines=true,
    captionpos=b,
    keepspaces=true,
    numbers=left,
    numbersep=5pt,
    showspaces=false,
    showstringspaces=false,
    showtabs=false,
    tabsize=2,
    frame=single,
    rulecolor=\color{codegray},
    language=bash,
    inputencoding=utf8,
    extendedchars=true,
    literate=
      {ä}{{\"a}}1
      {ö}{{\"o}}1
      {ü}{{\"u}}1
      {Ä}{{\"A}}1
      {Ö}{{\"O}}1
      {Ü}{{\"U}}1
      {ß}{{\ss}}1
}

% Header und Footer
\setlength{\headheight}{15pt}
\pagestyle{fancy}
\fancyhf{}
\rhead{Network-Klausur}
\lhead{\rightmark}
\cfoot{\thepage}

% Hyperref-Setup
\hypersetup{
    colorlinks=true,
    linkcolor=blue,
    filecolor=magenta,
    urlcolor=cyan,
}

\title{\textbf{Network-Vorlesung: Klausur-Vorbereitung}\\
Zusammenfassung Fakten, Wissensfragen und Praktische Aufgaben}
\author{Basierend auf den Vorlesungsunterlagen (sk1-sk10)}
\date{\today}

\begin{document}

\maketitle

\tableofcontents
\newpage

\part{Wichtige Fakten und Wissensfragen}

\chapter{Einführung in Netzwerke}

\section{Schlüsselbegriffe}
\begin{itemize}
    \item \textbf{Standards}: ISO/OSI, IEEE 802
    \item \textbf{Datennetze}: LAN, WAN, MAN
    \item \textbf{Switching}: VLAN, Port-Switching
    \item \textbf{Routing}: IP-Routing, RIP, OSPF
    \item \textbf{TCP/IP}: Protokoll-Familie für Datenkommunikation
\end{itemize}

\section{Wissensfragen}
\begin{enumerate}
    \item Was ist der Unterschied zwischen LAN und WAN?
    \item Erkläre das ISO/OSI-Referenzmodell (7 Schichten)
    \item Welche Standards/Normen sind für Netzwerke wichtig?
    \item Was versteht man unter Netzwerk-Adressierung?
    \item Welche Aufgaben hat eine Gateway?
\end{enumerate}

\chapter{Kommunikationsmodelle}

\section{Schlüsselbegriffe}
\begin{itemize}
    \item \textbf{Modulation}: Umwandlung von Signalen
    \item \textbf{Digitale Übertragung}: Binäre Kodierung
    \item \textbf{Duplex-Betrieb}: Halb- und Vollduplex
    \item \textbf{Multiplex}: Zeitmultiplex (TDM), Frequenzmultiplex (FDM)
    \item \textbf{BER}: Bitfehlerrate, Qualitätsmass
\end{itemize}

\section{Wissensfragen}
\begin{enumerate}
    \item Was ist der Unterschied zwischen analogen und digitalen Signalen?
    \item Erkläre Modulation und Demodulation
    \item Wann nutzt man Halb-Duplex vs.\ Vollduplex?
    \item Wie funktioniert Multiplexing?
    \item Wie wird Bitfehlerrate (BER) gemessen und bewertet?
\end{enumerate}

\chapter{Verkabelung und Übertragungsmedien}

\section{Schlüsselbegriffe}
\begin{itemize}
    \item \textbf{Kupferkabel}: Twisted Pair, Koaxial
    \item \textbf{LWL (Lichtwellenleiter)}: Glasfaser, Multimode/Singlemode
    \item \textbf{Wireless}: Funkübertragung, Mikrowelle
    \item \textbf{Strukturierte Verkabelung}: TIA/EIA-568
    \item \textbf{Dämpfung}: Signalverlust, Reichweite
\end{itemize}

\section{Wissensfragen}
\begin{enumerate}
    \item Vergleiche Kupferkabel, Koaxial und Glasfaser (Vor-/Nachteile)
    \item Warum nutzt man geschirmte vs.\ ungeschirmte Kabel?
    \item Wie funktioniert Lichtwellenleitung?
    \item Was ist das ``Cabling Standard'' TIA/EIA-568?
    \item Berechne Reichweite und Dämpfung bei verschiedenen Medien
\end{enumerate}

\chapter{Netzwerk-Interfaces}

\section{Schlüsselbegriffe}
\begin{itemize}
    \item \textbf{NIC}: Netzwerk-Schnittstellenkarte
    \item \textbf{Ethernet}: 10Base-T, 100Base-TX, Gigabit
    \item \textbf{MAC-Adresse}: Schicht-2-Adressierung
    \item \textbf{Switching}: MAC-Tabellen, Port-Learning
    \item \textbf{PHY/MAC-Layer}: Physikalische und Sicherungsschicht
\end{itemize}

\section{Wissensfragen}
\begin{enumerate}
    \item Erkläre den Aufbau und die Funktion einer MAC-Adresse
    \item Wie funktioniert das MAC-Learning in Switches?
    \item Vergleiche unterschiedliche Ethernet-Standards
    \item Was ist eine VLAN und wozu wird es genutzt?
    \item Wie wird bei Switches das Spanning Tree Protokoll verwendet?
\end{enumerate}

\chapter{Local Area Networks (LAN)}

\section{Schlüsselbegriffe}
\begin{itemize}
    \item \textbf{Switching}: Layer-2-Schicht, MAC-Bridging
    \item \textbf{VLAN}: Virtual LAN, Segmentierung
    \item \textbf{Spanning Tree}: Loop-Vermeidung (STP, RSTP)
    \item \textbf{VTP}: VLAN Trunking Protocol
    \item \textbf{Link Aggregation}: Port-Bündelung (LAG, 802.3ad)
\end{itemize}

\section{Wiskensfragen}
\begin{enumerate}
    \item Wie unterscheiden sich Hubs, Bridges und Switches?
    \item Erkläre VLAN-Konzept und 802.1Q-Tagging
    \item Warum brauchen wir Spanning Tree Protocol?
    \item Wie funktioniert RSTP (Rapid Spanning Tree)?
    \item Was ist Port-Bündelung (Link Aggregation) und wofür?
\end{enumerate}

\chapter{Internetworking und Routing}

\section{Schlüsselbegriffe}
\begin{itemize}
    \item \textbf{IP-Adressen}: IPv4, IPv6, Subnetting
    \item \textbf{Routing-Protokolle}: RIP, OSPF, BGP
    \item \textbf{Default Gateway}: Router-Funktion
    \item \textbf{Routing-Tabellen}: Nächster Hop-Bestimmung
    \item \textbf{ARP}: Address Resolution Protocol
\end{itemize}

\section{Wissensfragen}
\begin{enumerate}
    \item Erkläre IPv4-Subnetting (CIDR-Notation)
    \item Wie funktioniert ARP (Address Resolution Protocol)?
    \item Vergleiche Distanz-Vektor vs.\ Link-State Routing
    \item Wie arbeitet OSPF vs.\ BGP?
    \item Berechne Subnetz-Adressen und Broadcast-Adressen
\end{enumerate}

\chapter{TCP/IP-Protokoll}

\section{Schlüsselbegriffe}
\begin{itemize}
    \item \textbf{TCP}: Transmission Control Protocol (verbindungsorientiert)
    \item \textbf{UDP}: User Datagram Protocol (verbindungslos)
    \item \textbf{IP}: Internet Protocol, Routing-Protokoll
    \item \textbf{DNS}: Domain Name System
    \item \textbf{DHCP}: Dynamic Host Configuration Protocol
    \item \textbf{NAT}: Network Address Translation
\end{itemize}

\section{Wissensfragen}
\begin{enumerate}
    \item Unterschied zwischen TCP und UDP?
    \item Wie funktioniert ein TCP-Three-Way-Handshake?
    \item Erkläre die TCP-Zustandsübergänge
    \item Wie arbeitet DHCP (Discover, Offer, Request, Acknowledge)?
    \item Was ist NAT und wofür wird es verwendet?
\end{enumerate}

\chapter{IPv6 und erweiterte TCP/IP}

\section{Schlüsselbegriffe}
\begin{itemize}
    \item \textbf{IPv6}: Nächste Generation IP
    \item \textbf{IPv6-Adressen}: 128-Bit, Präfixlänge /64
    \item \textbf{Auto-Configuration}: SLAAC, DHCPv6
    \item \textbf{Neighbor Discovery}: ICMPv6
    \item \textbf{Dual Stack}: IPv4 und IPv6 parallel
\end{itemize}

\section{Wissensfragen}
\begin{enumerate}
    \item Warum wurde IPv6 entwickelt?
    \item Erkläre IPv6-Adressformat und Präfixlänge
    \item Wie funktioniert SLAAC (Stateless Address Autoconfiguration)?
    \item Unterschied DHCPv6 und SLAAC?
    \item Wie migriert man von IPv4 zu IPv6?
\end{enumerate}

\chapter{Multicast und Broadcast}

\section{Schlüsselbegriffe}
\begin{itemize}
    \item \textbf{Broadcast}: 1-zu-viele Kommunikation (LAN-Ebene)
    \item \textbf{Multicast}: gezielter 1-zu-viele (Gruppe)
    \item \textbf{IGMP}: Internet Group Management Protocol
    \item \textbf{TTL}: Time-to-Live (Multicast-Bereichs-Begrenzung)
    \item \textbf{Multicast-Adressen}: 224.0.0.0 - 239.255.255.255 (IPv4)
\end{itemize}

\section{Wissensfragen}
\begin{enumerate}
    \item Unterschied Broadcast und Multicast?
    \item Wie arbeitet IGMP?
    \item Welche Probleme entstehen durch Broadcast?
    \item Wie begrenzt man Multicast-Reichweite mit TTL?
    \item Gib Beispiele für Multicast-Anwendungen
\end{enumerate}

\chapter{Mobilnetze und Wireless}

\section{Schlüsselbegriffe}
\begin{itemize}
    \item \textbf{WLAN}: 802.11 (a/b/g/n/ac/ax)
    \item \textbf{Funkkanäle}: Frequenzbänder, Kanalbelegung
    \item \textbf{SSID}: Service Set Identifier
    \item \textbf{Authentifizierung}: WPA2, WPA3
    \item \textbf{Roaming}: Handoff zwischen APs
\end{itemize}

\section{Wissensfragen}
\begin{enumerate}
    \item Erkläre 802.11-Standards und deren Unterschiede
    \item Wie funktioniert WLAN-Authentifizierung (WPA2/WPA3)?
    \item Welche Kanäle und Frequenzen werden genutzt?
    \item Was ist der Unterschied zwischen SSID-Hidden und Open?
    \item Wie funktioniert Roaming in WLAN-Netzwerken?
\end{enumerate}

\newpage
\part{Praktische Aufgabentypem für die Klausur}

\chapter{Aufgabentyp 1: Netzwerk-Design und -Planung}

\section{Szenario 1.1: LAN für Bürogebäude}

Ein Bürogebäude mit 250 Mitarbeitern benötigt ein neues LAN-System:
\begin{itemize}
    \item 5 Stockwerke à 50 Mitarbeiter
    \item Zentrale IT-Infrastruktur in Etage 3
    \item Verschiedene Abteilungen (Verwaltung, Entwicklung, Support)
    \item Anforderung: Hohe Verfügbarkeit und Geschwindigkeit
\end{itemize}

\subsection{Aufgaben}
\begin{enumerate}
    \item Entwerfe eine Netzwerk-Topologie (welche Topologie auf welcher Schicht?)
    \item Wähle geeignete Hardware (Switches, Router, Cabling) und begründe
    \item Dimensioniere die erforderliche Bandbreite pro Etage
    \item Erstelle ein VLAN-Schema (Trennung nach Abteilungen)
    \item Plane Redundanz und Backup-Pfade
    \item Berechne Netzwerk-Adressen und Subnetze
\end{enumerate}

\subsection{Lösungsansätze}
\begin{itemize}
    \item Hierarchisches Design: Core-Distribution-Access
    \item Glasfaser zwischen Etagen, Cat6A in Büros
    \item PoE für Access Points
    \item VLAN pro Abteilung (Admin, Dev, Support, Guest)
    \item Redundante Uplinks, STP/RSTP aktivieren
\end{itemize}

\section{Szenario 1.2: Campus-Netzwerk mit mehreren Gebäuden}

Ein Unternehmen mit 3 Gebäuden (200m Entfernung) benötigt ein verbundenes Netzwerk:
\begin{itemize}
    \item Gebäude A: Verwaltung (100 Mitarbeiter)
    \item Gebäude B: Produktion (150 Mitarbeiter)
    \item Gebäude C: Lager (30 Mitarbeiter)
    \item Zentrale Services in Gebäude A
\end{itemize}

\subsection{Aufgaben}
\begin{enumerate}
    \item Entwerfe eine Multi-Building Netzwerk-Architektur
    \item Wie verbindest du die Gebäude? (Kupfer, Glasfaser, Funk?)
    \item Welche Redundanz-Strategien sind nötig?
    \item Wie segmentierst du die Abteilungen übergreifend?
    \item Wie migrierst du bestehende Netzwerke?
\end{enumerate}

\subsection{Lösungsansätze}
\begin{itemize}
    \item Glasfaser zwischen Gebäuden (Zukunftssicher)
    \item Geschirmte Cat6/Cat6A innerhalb Gebäude
    \item Redundante Uplinks zwischen Gebäuden
    \item VLANs mit Inter-VLAN-Routing (Router)
    \item Load-Balancing auf Kern-Switches
\end{itemize}

\chapter{Aufgabentyp 2: Protokoll-Analyse und Packet-Tracing}

\section{Szenario 2.1: TCP-Verbindung analysieren}

Gegeben ist eine TCP-Verbindung zwischen:
\begin{lstlisting}
Client: 192.168.1.50:54321
Server: 8.8.8.8:443
\end{lstlisting}

\subsection{Aufgaben}
\begin{enumerate}
    \item Erkläre den Three-Way-Handshake (SYN, SYN-ACK, ACK)
    \item Zeichne die Zustände des TCP-State-Machines
    \item Was passiert bei Paketverlusten (Retransmission)?
    \item Wie werden ACK-Nummern berechnet?
    \item Beschreibe die Trennung der Verbindung (4-Way-Handshake)
\end{enumerate}

\subsection{Erwartete Antworten}
\begin{itemize}
    \item SYN mit Initial Sequence Number (ISN)
    \item SYN-ACK mit Bestätigung und eigenem ISN
    \item ACK-Nummer = empfangene Sequence + Datenlänge
    \item TCP-Fenster (Window Size) für Durchsatzoptimierung
\end{itemize}

\section{Szenario 2.2: UDP vs.\ TCP für DNS}

Ein DNS-Client möchte eine Domain auflösen:
\begin{lstlisting}
Query: example.com -> IP-Adresse
Client: 192.168.1.50:12345
DNS-Server: 8.8.8.8:53
\end{lstlisting}

\subsection{Aufgaben}
\begin{enumerate}
    \item Warum verwendet DNS Port 53/UDP (statt TCP)?
    \item Wenn die Response > 512 Bytes ist, was passiert dann?
    \item Zeichne das DNS-Nachrichtenformat
    \item Wie funktioniert DNS-Caching?
    \item Was ist der Unterschied zwischen rekursiven und iterativen Queries?
\end{enumerate}

\subsection{Erwartete Antworten}
\begin{itemize}
    \item UDP für schnelle Abfragen, TCP für Zone-Transfers
    \item Truncation (TC-Bit) in UDP-Response
    \item DNS nutzt zusätzlich TCP für größere Responses
    \item TTL (Time-to-Live) für Cache-Validierung
\end{itemize}

\chapter{Aufgabentyp 3: Netzwerk-Fehlerbehebung (Troubleshooting)}

\section{Szenario 3.1: Konnektivitätsprobleme in mehreren Etagen}

Problem: Benutzer im 2.\ Stock können Server im 3.\ Stock nicht erreichen

\begin{lstlisting}
Client (2. Stock):   192.168.2.10 (Subnet: /24)
Server (3. Stock):   192.168.3.10 (Subnet: /24)
Router-2:            192.168.2.1
Router-3:            192.168.3.1
\end{lstlisting}

\subsection{Aufgaben}
\begin{enumerate}
    \item Beschreibe dein Troubleshooting-Vorgehen (Schritt-für-Schritt)
    \item Welche Kommandos würdest du einsetzen?
    \begin{itemize}
        \item Ping, Tracert/Traceroute, Pathping
        \item ipconfig/ifconfig, netstat
        \item arp -a, route print
    \end{itemize}
    \item Was sind mögliche Fehlerquellen auf verschiedenen Schichten?
    \begin{itemize}
        \item Layer 1: Kabel, LED-Status
        \item Layer 2: VLAN-Tagging, Spanning Tree
        \item Layer 3: Routing-Konfiguration, Firewall-Regeln
    \end{itemize}
    \item Wie überprüfst du Switch-Konfiguration?
\end{enumerate}

\subsection{Erwartete Antworten}
\begin{enumerate}
    \item Lokale Konnektivität prüfen (ping 192.168.2.1)
    \item Default-Gateway überprüfen
    \item Routing-Tabelle analysieren (route print)
    \item Auf Router: Routing zwischen den Subnetzen?
    \item VLAN-Tagging auf trunk ports richtig?
    \item ACL oder Firewall-Regeln blockieren?
\end{enumerate}

\section{Szenario 3.2: DHCP-Ausfallszenario}

Problem: Nach Neustart des DHCP-Servers können neue Clients keine IP-Adresse erhalten

\subsection{Aufgaben}
\begin{enumerate}
    \item Wie überprüfst du DHCP-Funktionalität?
    \item Was sind potenzielle Probleme?
    \begin{itemize}
        \item DHCP-Server läuft nicht
        \item DHCP-Relay ist nicht konfiguriert
        \item DHCP-Spool ist leer
        \item Broadcast-Blockierung
    \end{itemize}
    \item Wie debuggst du DHCP-Kommunikation?
    \item Wie konfigurierst du DHCP-Failover/Redundanz?
\end{enumerate}

\subsection{Erwartete Antworten}
\begin{itemize}
    \item DHCP-Pakete mit Wireshark/TCPdump analysieren
    \item DHCP-Server-Logs prüfen
    \item DHCP-Relay auf Routern überprüfen
    \item Auf Client: \texttt{ipconfig /renew} (Windows) oder \texttt{dhclient} (Linux)
    \item Reservierte IP-Bereiche für DHCP definieren
\end{itemize}

\chapter{Aufgabentyp 4: Performance-Analyse und Optimierung}

\section{Szenario 4.1: Netzwerk wird immer langsamer}

Ein Administrator misst folgende Metriken:
\begin{lstlisting}
Verfügbare Bandbreite:    1 Gbps
Aktuelle Auslastung:      85%
Paket-Verluste:           2%
Latenz:                   150ms
Jitter:                   50ms
\end{lstlisting}

\subsection{Aufgaben}
\begin{enumerate}
    \item Sind diese Werte problematisch? Worauf kommt es an?
    \item Was sind mögliche Ursachen für hohe Latenz?
    \item Welche Optimierungsmaßnahmen schlägst du vor?
    \item Wie implementierst du QoS (Quality of Service)?
    \item Berechne Durchsatz bei verschiedenen Verlustraten
\end{enumerate}

\subsection{Erwartete Antworten}
\begin{itemize}
    \item 85\% Auslastung ist bedenklich (Puffering)
    \item 2\% Packet Loss ist zu hoch (sollte < 0.1\%)
    \item 150ms Latenz ist problematisch für VoIP/Video
    \item QoS: Prioritäten setzen (Voice, Video, Data)
    \item Traffic-Shaping, Rate-Limiting für nicht-kritische Dienste
\end{itemize}

\section{Szenario 4.2: WLAN-Performance-Problem}

In einem Büro nutzen 50 Benutzer ein WLAN-Netzwerk, viele haben Verbindungsprobleme

\subsection{Aufgaben}
\begin{enumerate}
    \item Was sind typische Ursachen für schlechtes WLAN?
    \item Wie misst man WLAN-Signal-Qualität?
    \begin{itemize}
        \item RSSI (Received Signal Strength Indicator)
        \item SNR (Signal-to-Noise Ratio)
        \item Data Rate
    \end{itemize}
    \item Wie optimierst du Kanal-Belegung (2,4 GHz vs.\ 5 GHz)?
    \item Wie viele APs brauchst du und wo platzierst du sie?
    \item Welche 802.11-Standards wählst du?
\end{enumerate}

\subsection{Erwartete Antworten}
\begin{itemize}
    \item Zu viele Benutzer pro AP $\rightarrow$ zusätzliche APs
    \item Interferenzen mit anderen WLANs $\rightarrow$ Kanal-Wechsel
    \item Schwaches Signal $\rightarrow$ Antenne/Position anpassen
    \item 5 GHz hat höhere Bandbreite, geringere Reichweite
    \item 802.11ax (Wi-Fi 6) für hohe Benutzeranzahl
\end{itemize}

\chapter{Aufgabentyp 5: Netzwerk-Sicherheit und Segmentierung}

\section{Szenario 5.1: Sichere Netzwerk-Architektur}

Ein Unternehmen möchte sein Netzwerk besser sichern. Anforderungen:
\begin{itemize}
    \item Trennung von Gast- und Unternehmens-Netzwerk
    \item Schutz vor DDoS-Attacken
    \item Sichere Fernzugriffe für Mitarbeiter (VPN)
    \item Einhaltung von Compliance-Anforderungen
\end{itemize}

\subsection{Aufgaben}
\begin{enumerate}
    \item Entwerfe eine DMZ-Architektur (3-Schichten-Modell)
    \item Wie funktioniert VLAN-Segmentierung?
    \begin{itemize}
        \item Management-VLAN
        \item Guest-VLAN
        \item Production-VLAN
        \item Sensitive-Data-VLAN
    \end{itemize}
    \item Konfiguriere Inter-VLAN-Routing mit Access Control Lists (ACL)
    \item Wie implementierst du VPN (Remote Access)?
    \begin{itemize}
        \item Welche VPN-Technologien (IPSec, SSL/TLS)?
        \item Authentifizierung (Zertifikate, 2FA)?
    \end{itemize}
    \item DDoS-Schutzmechanismen:
    \begin{itemize}
        \item Ratelimiting
        \item IP-Reputation-Filtering
        \item Geographische Filterung
    \end{itemize}
\end{enumerate}

\subsection{Erwartete Antworten}
\begin{itemize}
    \item DMZ: Public Services (Web, Mail)
    \item Internal Network: Unternehmenssysteme
    \item Management Network: Admin-Zugriff nur
    \item ACL-Beispiel: Gast-VLAN darf nicht auf Production
\end{itemize}

\section{Szenario 5.2: Wireless-Sicherheit}

Sichere die WLAN-Infrastruktur gegen bekannte Angriffe

\subsection{Aufgaben}
\begin{enumerate}
    \item Erkläre WPA2 vs.\ WPA3 Unterschiede
    \item Wie konfigurierst du starke Authentifizierung?
    \begin{itemize}
        \item Pre-Shared Key (PSK) vs.\ Enterprise (802.1X)
    \end{itemize}
    \item Was ist TKIP und warum ist es obsolet?
    \item Wie funktioniert 802.1X mit RADIUS?
    \item Welche weiteren Sicherheitsmaßnahmen?
    \begin{itemize}
        \item MAC-Filtering
        \item SSID-Cloaking (Sicherheit durch Obscurity)
        \item Encryption: AES-CCMP
    \end{itemize}
\end{enumerate}

\subsection{Erwartete Antworten}
\begin{itemize}
    \item WPA3: Individualized Data Encryption (IDE), Protection Management Frames
    \item Enterprise WPA2 mit RADIUS für Unternehmen
    \item PSK OK für kleine Netzwerke, Enterprise für Sicherheit
\end{itemize}

\chapter{Aufgabentyp 6: IP-Adressierung und Subnetting}

\section{Szenario 6.1: Subnetting für Multi-Site-Netzwerk}

Ein Unternehmen hat die IP-Range: \textbf{10.0.0.0/16}

Sie benötigt:
\begin{itemize}
    \item Verwaltung: 50 Hosts
    \item Produktion: 200 Hosts
    \item Entwicklung: 80 Hosts
    \item Gast: 30 Hosts
\end{itemize}

\subsection{Aufgaben}
\begin{enumerate}
    \item Entwerfe ein Subnetting-Schema (VLSM -- Variable Length Subnet Masking)
    \item Berechne für jedes Subnet:
    \begin{itemize}
        \item Netzwerk-Adresse
        \item Broadcast-Adresse
        \item Erste/letzte verwendbare Host-Adresse
        \item Subnet-Mask (CIDR-Notation)
    \end{itemize}
    \item Wieviele Bits brauchst du für Hosts bei größtem Subnet?
    \item Wieviele verfügbare Subnetze hast du noch?
\end{enumerate}

\subsection{Erwartete Antworten}
\begin{lstlisting}
Verwaltung (50 Hosts):    10.0.0.0/26
  (Netzwerk: 10.0.0.0, Broadcast: 10.0.0.63)
Produktion (200 Hosts):   10.0.1.0/23
  (Netzwerk: 10.0.1.0, Broadcast: 10.0.2.255)
Entwicklung (80 Hosts):   10.0.3.0/25
  (Netzwerk: 10.0.3.0, Broadcast: 10.0.3.127)
Gast (30 Hosts):          10.0.4.0/27
  (Netzwerk: 10.0.4.0, Broadcast: 10.0.4.31)
\end{lstlisting}

\section{Szenario 6.2: IPv6-Migration}

Plane die Migration von IPv4 zu IPv6

\subsection{Aufgaben}
\begin{enumerate}
    \item Dein Unternehmen hat die IPv6-Range: \textbf{2001:db8::/32}
    \item Design ein Addressing-Schema:
    \begin{itemize}
        \item Site ID für verschiedene Standorte
        \item Subnet ID für verschiedene Funktionen
        \item Host ID (meist auto-konfiguriert)
    \end{itemize}
    \item Berechne Adressen für:
    \begin{itemize}
        \item Verwaltungs-Subnet (Berlin)
        \item Produktion-Subnet (München)
        \item Guest-Subnet (beide Standorte)
    \end{itemize}
    \item Wie konfigurierst du Dual-Stack (IPv4+IPv6)?
    \item Wie funktioniert SLAAC?
\end{enumerate}

\subsection{Erwartete Antworten}
\begin{lstlisting}
IPv6-Struktur: 2001:db8:site:subnet:interface:host
Berlin Verwaltung: 2001:db8:0001:1::/64
München Produktion: 2001:db8:0002:2::/64
\end{lstlisting}

\chapter{Aufgabentyp 7: Switch- und Router-Konfiguration}

\section{Szenario 7.1: VLAN- und Trunk-Konfiguration}

Konfiguriere folgende Netzwerk-Infrastruktur:

\begin{center}
\begin{tabular}{cccccc}
Core Switch & & & & & \\
/ & | & $\backslash$ & & & \\
{[}SW1{]} & {[}SW2{]} & {[}SW3{]} & & & \\
\multicolumn{6}{c}{{[}Clients auf verschiedenen Etagen{]}} \\
\end{tabular}
\end{center}

\subsection{Aufgaben}
\begin{enumerate}
    \item Erstelle VLANs:
    \begin{itemize}
        \item VLAN 10: Verwaltung (Netzwerk: 192.168.10.0/24)
        \item VLAN 20: Produktion (Netzwerk: 192.168.20.0/24)
        \item VLAN 30: Entwicklung (Netzwerk: 192.168.30.0/24)
        \item VLAN 99: Management (Netzwerk: 192.168.99.0/24)
    \end{itemize}
    \item Konfiguriere Trunk-Links zwischen Switches (802.1Q)
    \item Konfiguriere Access-Ports für Clients (welches VLAN?)
    \item Aktiviere Spanning Tree Protocol (STP)
    \begin{itemize}
        \item Wähle Root Bridge
        \item Überprüfe Loop-Freiheit
    \end{itemize}
    \item Inter-VLAN-Routing:
    \begin{itemize}
        \item Welches Gerät? (Router oder Layer-3-Switch)
        \item Konfiguriere Sub-Interfaces auf Router
    \end{itemize}
\end{enumerate}

\subsection{Erwartete Antworten (vereinfacht)}
\begin{lstlisting}[language=bash]
[Core Switch - Cisco IOS Beispiel]
vlan 10
name Verwaltung
vlan 20
name Produktion

[Trunk zu anderen Switches]
interface GigabitEthernet0/1
switchport mode trunk
switchport trunk allowed vlan 10,20,30,99

[Access Port für Client]
interface FastEthernet0/10
switchport mode access
switchport access vlan 10
\end{lstlisting}

\section{Szenario 7.2: Routing-Konfiguration}

Verbinde mehrere Subnetze mit Routing:

\begin{center}
\begin{tabular}{ccc}
Verwaltung & & Produktion \\
192.168.10.0/24 & & 192.168.20.0/24 \\
$|$ & & $|$ \\
{[}Router A{]} & {-}{-}{-} & {[}Router B{]} \\
(HQ) & & (Filiale)
\end{tabular}
\end{center}

\subsection{Aufgaben}
\begin{enumerate}
    \item Konfiguriere Router-Schnittstellen (IP-Adressen)
    \item Wähle Routing-Protokoll:
    \begin{itemize}
        \item RIP (einfach, aber langsam)
        \item OSPF (dynamisch, komplexer)
        \item BGP (für Inter-AS-Routing)
    \end{itemize}
    \item Konfiguriere das Protokoll (OSPF Beispiel):
    \begin{itemize}
        \item Area Definition
        \item Interface-Konfiguration
        \item Metric-Adjustment
    \end{itemize}
    \item Überprüfe Routing-Tabellen
    \item Teste Konnektivität mit Ping/Tracert
\end{enumerate}

\subsection{Erwartete Antworten (vereinfacht)}
\begin{lstlisting}[language=bash]
[Router A]
router ospf 1
network 192.168.10.0 0.0.0.255 area 0
network 10.0.0.0 0.0.0.255 area 0

[Router B]
router ospf 1
network 192.168.20.0 0.0.0.255 area 0
network 10.0.0.0 0.0.0.255 area 0
\end{lstlisting}

\chapter{Aufgabentyp 8: Komplexe Szenario-Aufgaben}

\section{Szenario 8.1: Hybrid-Netzwerk mit Cloud-Integration}

Ein Unternehmen benötigt ein Netzwerk mit Cloud-Anbindung:

\subsection{Anforderungen}
\begin{itemize}
    \item Lokales Netzwerk: 200 Benutzer
    \item Cloud-Services (AWS/Azure): Geschäftskritische Anwendungen
    \item Remote-Mitarbeiter: Sichere VPN-Verbindung
    \item Hohe Verfügbarkeit: Redundante Pfade
\end{itemize}

\subsection{Aufgaben}
\begin{enumerate}
    \item Entwerfe die Gesamtarchitektur:
    \begin{itemize}
        \item On-Premise Netzwerk
        \item Cloud-Verbindung (DCI/Carrier?)
        \item Remote-Access (VPN)
        \item Management-Netzwerk
    \end{itemize}
    \item Netzwerk-Design:
    \begin{itemize}
        \item Welche Topologie (Mesh, Hub-Spoke)?
        \item Redundanze-Anforderungen?
        \item WAN-Technologien (MPLS, SD-WAN)?
    \end{itemize}
    \item Sicherheit:
    \begin{itemize}
        \item Firewall-Positionierung
        \item VPN-Protokolle und Authentifizierung
        \item Cloud-Security-Anforderungen
        \item Zero-Trust-Model?
    \end{itemize}
    \item Performance:
    \begin{itemize}
        \item Bandbreiten-Dimensionierung
        \item QoS-Anforderungen
        \item Cloud-Latenz-Anforderungen
        \item Caching-Strategien
    \end{itemize}
    \item Compliance:
    \begin{itemize}
        \item Datenschutz (DSGVO)
        \item Audit-Anforderungen
        \item Backup/Disaster-Recovery
    \end{itemize}
\end{enumerate}

\subsection{Erwartete Struktur}
\begin{lstlisting}
[On-Premise]          [Firewall]        [Cloud Provider]
- Clients        ->    - Security        - Applications
- Server             - VPN-Gateway      - Storage
- Storage            - Load Balancer    - Database

[Remote Users] --VPN--> [Firewall] --Tunnel--> [Cloud]
\end{lstlisting}

\section{Szenario 8.2: Optimierung eines langsamen Netzwerks}

Ein bestehendes Netzwerk läuft ineffizient:

\subsection{Aktuelle Situation}
\begin{lstlisting}
- 10 Mbps WAN-Verbindung zwischen 2 Standorten
- 85% Auslastung
- Video-Konferenz Probleme
- Backup dauert zu lange
\end{lstlisting}

\subsection{Aufgaben}
\begin{enumerate}
    \item Analysiere die Probleme:
    \begin{itemize}
        \item Was ist der Bottleneck?
        \item Welche Services verursachen hohe Last?
        \item Sind Redundanzen vorhanden?
    \end{itemize}
    \item Konfiguriere QoS:
    \begin{itemize}
        \item Priorisierung (Video > Voice > Data)
        \item Rate-Limiting für Backup
        \item Burst-Behandlung
    \end{itemize}
    \item Optimierungsmaßnahmen:
    \begin{itemize}
        \item WAN-Verbesserung (Bandbreite upgraden?)
        \item Traffic-Lokalisierung (Caching, Local Storage)
        \item Compression und Deduplication
        \item SD-WAN Evaluierung
    \end{itemize}
    \item Berechne neue Anforderungen:
    \begin{itemize}
        \item Bandbreite für Video-Konferenz (HD, 20 Nutzer)
        \item Backup-Anforderungen (RPO/RTO)
        \item Redundanz-Anforderungen
    \end{itemize}
    \item Implementierungs-Plan:
    \begin{itemize}
        \item Phased Rollout
        \item Testing und Validierung
        \item Monitoring und Alerting
    \end{itemize}
\end{enumerate}

\chapter{Aufgabentyp 9: Fehleranalyse und Systemausfallszenario}

\section{Szenario 9.1: Kompletter Netzwerk-Ausfall}

Das gesamte Netzwerk fällt aus. Wie gehst du vor?

\subsection{Aufgaben}
\begin{enumerate}
    \item Sofortmaßnahmen (erste 30 Minuten):
    \begin{itemize}
        \item Wie priorisierst du Recovery?
        \item Was überprüfst du zuerst?
        \item Kommunikation mit Nutzern?
    \end{itemize}
    \item Detaillierte Fehlersuche:
    \begin{itemize}
        \item Schicht-für-Schicht Diagnose
        \item Welche Tools/Kommandos?
        \item Monitoring-Daten analysieren
    \end{itemize}
    \item Fehlerszenarien und Lösungen:
    \begin{itemize}
        \item \textbf{Szenario 1}: Kern-Switch ausgefallen
        \item \textbf{Szenario 2}: Internet-Gateway fehlt
        \item \textbf{Szenario 3}: DHCP-Server nicht erreichbar
        \item \textbf{Szenario 4}: DNS-Ausfall
    \end{itemize}
    \item Post-Mortem:
    \begin{itemize}
        \item Ursachenanalyse (RCA)
        \item Wie hätte man das verhindern können?
        \item Redundanz-Verbesserungen
    \end{itemize}
\end{enumerate}

\newpage
\part{Wichtige Formeln und Berechnungen}

\chapter{Networking-Formeln}

\section{Durchsatz-Berechnung}
\[
\text{Effektiver Durchsatz} = \text{Bandbreite} \times (1 - \text{Packet-Loss-Rate})
\]

\textit{Beispiel:} $100 \text{ Mbps} \times (1 - 0.02) = 98 \text{ Mbps}$

\section{RTT (Round Trip Time) / Latenz}
\[
\text{RTT} = 2 \times \frac{\text{Distanz}}{\text{Lichtgeschwindigkeit}}
\]

\textit{Beispiel:} $100 \text{ km} \rightarrow \text{RTT} \approx 2 \times \frac{100.000 \text{ m}}{300.000.000 \text{ m/s}} \approx 0.67 \text{ ms}$

\section{Bandbreite für Video-Streaming}
\[
\text{Bitrate (Mbps)} = \frac{\text{Auflösung} \times \text{FPS} \times \text{Bit-Tiefe}}{1.000.000}
\]

\textit{Beispiel:} HD (1280×720, 30 fps, 8-bit): $\approx 10$--$20$ Mbps

\section{TCP-Fenster-Größe und Durchsatz}
\[
\text{Durchsatz} = \frac{\text{Window Size}}{\text{RTT}}
\]

\textit{Beispiel:} $\frac{65.536 \text{ Bytes}}{50 \text{ ms}} = 1.3 \text{ Mbps}$

\section{Subnetting-Berechnung}
\[
\text{Verfügbare Hosts} = 2^n - 2
\]
wobei $n$ = Anzahl Host-Bits

\textit{Beispiel:} $/24 \rightarrow 8$ Host-Bits $\rightarrow 2^8 - 2 = 254$ Hosts

\part{Checklisten für Klausur-Vorbereitung}

\chapter{Verständnis-Checkliste}

\begin{itemize}
    \item[$\checkmark$] Kann ich das ISO/OSI-Modell erklären?
    \item[$\checkmark$] Kenne ich die wichtigsten Protokolle und deren Ports?
    \item[$\checkmark$] Kann ich Subnetze berechnen?
    \item[$\checkmark$] Verstehe ich den Unterschied zwischen Routing und Switching?
    \item[$\checkmark$] Kann ich VLAN-Konzept erklären?
    \item[$\checkmark$] Kenne ich die verschiedenen Netzwerk-Topologien?
    \item[$\checkmark$] Verstehe ich TCP vs.\ UDP?
    \item[$\checkmark$] Kann ich eine Netzwerk-Fehlersuche durchführen?
\end{itemize}

\chapter{Praktische-Fähigkeits-Checkliste}

\begin{itemize}
    \item[$\checkmark$] Kann ich eine Netzwerk-Architektur zeichnen?
    \item[$\checkmark$] Kann ich IP-Adressen berechnen?
    \item[$\checkmark$] Kann ich eine ACL schreiben?
    \item[$\checkmark$] Kann ich VLAN/Trunk konfigurieren?
    \item[$\checkmark$] Verstehe ich Routing-Tabellen?
    \item[$\checkmark$] Kann ich QoS konfigurieren?
    \item[$\checkmark$] Kann ich VPN aufsetzen?
    \item[$\checkmark$] Kann ich Fehlerbehebung durchführen?
\end{itemize}

\vfill
\begin{center}
\textbf{\Large Viel Erfolg bei der Klausur!}
\end{center}

\end{document}
